% ---------------------------------------------------------------------
%  TÓM TẮT
% ---------------------------------------------------------------------
\begin{abstract}

Việt Nam chịu ảnh hưởng trung bình nhiều cơn bão và áp thấp nhiệt đới mỗi
năm, nên biết sớm hướng đi của bão giúp chủ động phòng tránh. Đồ án xây dựng
một hệ thống trí tuệ nhân tạo gợi ý đường đi của cơn bão, phát biểu dưới dạng
bài toán dự báo chuỗi thời gian nhiều bước: từ bốn quan trắc gần nhất (24
giờ), hệ thống ước lượng vị trí tâm bão sau 6, 12, 24, 48 và 72 giờ kèm nón
bất định. Đồ án huấn luyện và so sánh bảy phương pháp, gồm hai mốc so sánh
(persistence, CLIPER), k láng giềng gần nhất dựa trên các cơn bão lịch sử
tương tự (k-NN analog), gradient boosting, LSTM, Attention-LSTM và
Transformer, trên 1.089 cơn bão ở Tây Bắc Thái Bình Dương giai đoạn
1980--2024, chia tập theo năm để không rò rỉ thông tin. Trên tập kiểm thử
2021--2024, các phương pháp học giảm sai số 72 giờ khoảng 200~km so với
persistence (676~km xuống 462--476~km), có ý nghĩa thống kê. Thí nghiệm cho
thấy học độ dịch chuyển thay vì tọa độ tuyệt đối giảm sai số 6 giờ từ 115~km
xuống 33~km, và bổ sung đặc trưng dòng dẫn đường khí quyển cho mô hình kết hợp
LSTM và Transformer giảm thêm 35~km tại 72 giờ. Hệ thống được triển khai thành
API và bản đồ web, mỗi gợi ý kèm danh sách cơn bão lịch sử tương tự để người
dùng hiểu vì sao có gợi ý đó. Kết quả cho thấy với dữ liệu chỉ có quỹ đạo, mô
hình học sâu không vượt trội phương pháp cổ điển; nút thắt nằm ở thông tin khí
quyển chứ không ở kiến trúc mạng.

\vspace{0.8em}
\noindent\textbf{Từ khoá:} dự báo quỹ đạo bão, chuỗi thời gian, LSTM,
Transformer, k láng giềng gần nhất, lập luận dựa trên tình huống.

\end{abstract}
